\scanpage{12}

% loc: construction
\uline{Contraction}\quad Soit $\Phi^c(C)$ la catégorie formée des \deleted{familles} couples $(\xi,\gamma)$, où $\xi\in\operatorname{Ob}\Phi(\unsure{C|1})$, $\xi=(I,(L_i)_{i\in I})$ (les $L_i\in\operatorname{Ob}C$), et où $\gamma$ est formé

% loc: data, part 1
1) d'un ensemble \inserted{$\gamma'$} (qui peut être vide) de parties de $I$, qui sont \unsure{chacune|de card.} de cardinal 2 et deux à deux disjointes

% loc: data, part 2
2) Pour $\forall J\in\gamma'$, d'un hom de

% loc: pairing
\[
\begin{aligned}\Pi(\xi|J)&\longrightarrow1_C\\
\rotatebox{90}{$=$}&\\
L_{j'}\otimes L_{j''}&\\[-2pt]\text{si }J=\{j',j''\}&
\end{aligned}
\]

% loc: left, beside pairing
% slanted
\margin{$\bigcup\gamma'=\bigcup_{\alpha\in\gamma'}J_\alpha$\\ \unsure{est appelé l'ens} \unsure{des|de}\\ \unsure{ou|des} indices contractés de \deleted{\unsure{$(\Phi^c(C))$}}\inserted{\unsure{$\xi$|$\gamma$}}}

% loc: bracket under pairing
[qui soit un accouplement parfait\\ i.e. une dualité]

% loc: morphisms and restriction
\deleted{\unsure{Les}} Les hom (qui sont des isom) de $\Phi^c(C)$ clairs. On trouve un foncteur évident \inserted{$\Phi^c(C)\to\Phi(C)$} (\unsure{une paire de foncteurs|en plus des foncteurs|des foncteurs} d'inclusion inc : $\operatorname{res}:\Phi^c(C)\longrightarrow\Phi(C)$

% loc: forgetting the pairs
(oubli des indices contractés). $\Phi^c(\unsure{C|1})$ est aussi une $\otimes$-catégorie ass. comm. \deleted{Ceci dit,} unitaire, et res est un $\otimes$ foncteur $\unsure{\mathrm{AUC}|\mathrm{ANC}}$.

% loc: last paragraph
Ceci dit, on trouve un hom canonique

% loc: contraction arrow
\[\Pi\circ\operatorname{inc}\xrightarrow[\text{Contraction}]{c}\Pi\circ\operatorname{res}\]

% loc: bottom
de $\otimes$-foncteurs AUC de $\Phi^c(C)$ dans $C$.

\noindent\hfil\rule{2.7em}{.4pt}\hfil

% archive sheet number: 11
